\begin{python0}
from solutions import *; clear()
\end{python0}
\sectionthree{Why Overloading?}

The reason why we have function (and operator) overloading is clear: Many functions perform similar operations and we want to remember fewer names.

This is also the case in Math. Not only is + used for addition of real numbers, it's also used for things like addition of functions, addition of vectors, addition of matrices, etc. So having function/operator overloading in a programming language allows us to design software to match the mathematical language/symbols that we are used to.

Even in daily use of the English language, we see the same word used similar but different situations: open a door, open a jar, open my wallet, open your mind, etc.

Overloading is a software technique for controlling software complexities.

(Some older programming languages do not have function overloading. Some don't even allow you to define your own operators. Also, there are programming languages where overloading is not allowed. In fact there are programming languages where the symbol for adding integers is different from the symbol for adding doubles.)

